\sectioncentered*{Заключение}
\addcontentsline{toc}{section}{Заключение}

В лабораторной работе разработано приложение с обычными и конфиденциальными записями, выполнен анализ защищённости и построена методика, применённая к этой системе.

\textbf{Семантическая составляющая.} Защищённость складывается из мер на уровнях ОС, приложения, сети и данных. Разделение сущностей и шифрование тел секретов на диске повышают устойчивость к чтению файла базы, но не заменяют защиту ключа и канала.

\textbf{Прагматическая составляющая.} Получены воспроизводимый стенд, осмотр хранилища и методика с уровнями риска, позволяющие повторять оценку и формулировать меры: вынос ключа, TLS, отдельный пользователь ОС.

\begin{thebibliography}{9}
	\addcontentsline{toc}{section}{Список использованных источников}

	\bibitem{howard}
	Ховард М., Лебланк Д. Защищённый код. \mdash  М. : Русская редакция, 2004. \mdash  704~с.

	\bibitem{howard24}
	Ховард М., Лебланк Д., Вьега Дж. 24 смертных греха компьютерной безопасности. \mdash  СПб. : Питер, 2010. \mdash  400~с.

	\bibitem{panasenko}
	Панасенко С.\,П. Алгоритмы шифрования. Специальный справочник. \mdash  СПб. : БХВ-Петербург, 2009. \mdash  576~с.

	\bibitem{cwe}
	CWE: Common Weakness Enumeration [Электронный ресурс]. \mdash  Режим доступа: \url{http://cwe.mitre.org/}. \mdash  Дата доступа: 12.09.2026.

	\bibitem{owasp}
	OWASP Vulnerability Scanning Tools [Электронный ресурс]. \mdash  Режим доступа: \url{https://owasp.org/www-community/Vulnerability_Scanning_Tools}. \mdash  Дата доступа: 12.09.2026.

	\bibitem{flask}
	Flask documentation [Электронный ресурс]. \mdash  Режим доступа: \url{https://flask.palletsprojects.com/}. \mdash  Дата доступа: 12.09.2026.
\end{thebibliography}
